\subsection{Criterios de Selección}

Los criterios de selección se definieron en función del objetivo metodológico del
trabajo. Como la investigación requería modificar la forma de almacenamiento del
vector de estado sin alterar la semántica general de la simulación, se priorizaron
aquellos factores que incidían de manera más directa en la viabilidad de esa
intervención y en la posibilidad de extenderla posteriormente. La Tabla~\ref{tab:criterios_sim} resume los criterios definidos con su descripción y ponderación.

\begin{table}[H]
    \centering
    \small
    \setlength{\tabcolsep}{4pt}
    \caption{Criterios usados para la selección del simulador cuántico}
    \label{tab:criterios_sim}
    \begin{tabular}{>{\raggedright\arraybackslash}p{1.5cm}>{\raggedright\arraybackslash}p{4.3cm}>{\raggedright\arraybackslash}p{0.42\linewidth}>{\centering\arraybackslash}p{1.2cm}}
        \hline
        \textbf{Código} & \textbf{Criterio} & \textbf{Descripción} & \textbf{Peso} \\ \hline
        C1 & Intervenibilidad del almacenamiento del vector de estado &
        Evalúa qué tan fácilmente puede reemplazarse, envolverse o rediseñarse la
        estructura que almacena las amplitudes complejas. Este criterio recibió la mayor
        ponderación porque el aporte técnico principal del trabajo consiste en modificar
        la forma de almacenamiento del estado cuántico. &
        0.45 \\
        C2 & Desacoplamiento entre lógica de simulación y backend de memoria &
        Valora en qué medida la lógica asociada a compuertas, evolución del circuito y
        operaciones sobre el estado puede mantenerse estable cuando se altera la forma
        de almacenamiento subyacente. Una arquitectura con mayor separación entre
        estas responsabilidades favorece una integración con impacto controlado sobre
        el resto del simulador. &
        0.35 \\
        C3 & Proyección de escalabilidad y alineación con trabajo futuro &
        Considera la presencia de mecanismos de paralelismo, distribución o aceleración
        útiles en extensiones posteriores del trabajo. Aunque este aspecto no fue el eje
        de la implementación actual, se incorporó para situar la selección en un contexto
        realista de simulación cuántica y de posibles desarrollos futuros. &
        0.20 \\ \hline
    \end{tabular}
\end{table}

La ponderación asignada refleja la prioridad de intervenir el almacenamiento del vector
de estado y mantener control sobre esa intervención. Por esta razón, los criterios de
modificabilidad arquitectónica recibieron mayor peso que la proyección de
escalabilidad general.
